% GENERATED FILE. Edit canonical lessons, then run npm run generate:books.
% canonical-source-hash: fnv1a64:e8a02ac59b52ebaf
% canonical-lessons: LA-C416-iustitia, LA-C416-aequitas, LA-C416-obtempero, LA-C416-coerceo, LA-C416-vindico

\chapter{Justice and Obedience}
\label{ch:la-justice-and-obedience}
\hlchaptermodality{416}

\begin{chapteropening}
\textbf{By the end of this chapter:} \emph{I can talk about justice and fairness, say that someone obeys the law or restrains people, and that a wrong was avenged.}

Complete the last lesson of chapter 416: I can talk about justice and fairness, say that someone obeys the law or restrains people, and that a wrong was avenged.
\end{chapteropening}

\section[iūstitia, iūstitiae]{iūstitia, iūstitiae — justice}

\label{lesson:LA-C416-iustitia}



\begin{quote}
Before the new one: say the Latin for a whip, then the Latin for custody.
\end{quote}

\subsection*{You'll want to know: iūstitia, iūstitiae}
\textbf{iūstitia, iūstitiae} — \textquotedblleft{}justice\textquotedblright{}.

\begin{quote}
Iūstitia fundāmentum cīvitātis est.
\end{quote}

\begin{quote}
\emph{Justice is the foundation of the state.}
\end{quote}

\emph{iūstitia, iūstitiae} f., built on \emph{iūstus}, fair and right. English \emph{justice}.

\begin{cousinweb}
\emph{iūstitia} is built on \emph{iūstus} \textquotedblleft{}righteous, fair\textquotedblright{} (EtymDB), which comes from \emph{iūs} \textquotedblleft{}right, law\textquotedblright{} (Lewis \& Short). English \emph{justice} comes from it through French.
\end{cousinweb}

\begin{grammarlens}[title={What you've built}]
One more word for telling, explaining and describing at B1.
\end{grammarlens}

\subsection*{Your turn}
\begin{itemize}
\raggedright
  \item \emph{Say it:} \emph{iūstitia, iūstitiae}
  \item \emph{Say it:} \emph{iūstitia, iūstitiae}, once more
  \item \emph{Say it:} a sentence of your own with \emph{iūstitia, iūstitiae}
  \item \emph{Recall:} say the Latin for a whip, then the Latin for custody, then say \emph{iūstitia, iūstitiae} again
\end{itemize}

\begin{tcolorbox}[breakable,colback=teal!4,colframe=teal!35!black,title={Before you move on}]
What does iūstitia, iūstitiae mean? (\textquotedblleft{}Justice\textquotedblright{}.) Say it once more.
\end{tcolorbox}
\section[aequitās, aequitātis]{aequitās, aequitātis — fairness, equity}

\label{lesson:LA-C416-aequitas}



\begin{quote}
Before the new one: say the Latin for custody, then the Latin for justice.
\end{quote}

\subsection*{You'll want to know: aequitās, aequitātis}
\textbf{aequitās, aequitātis} — \textquotedblleft{}fairness, equity\textquotedblright{}.

\begin{quote}
Iūdex magnā aequitāte iūdicāvit.
\end{quote}

\begin{quote}
\emph{The judge gave judgement with great fairness.}
\end{quote}

\emph{aequitās, aequitātis} f., built on \emph{aequus}, level, fair. English \emph{equity}.

\begin{cousinweb}
\emph{aequitās} is built on \emph{aequus} \textquotedblleft{}level, fair\textquotedblright{} (Lewis \& Short), with the ending \emph{-tās} that makes a quality. English \emph{equity} comes from it.
\end{cousinweb}

\begin{grammarlens}[title={What you've built}]
One more word for telling, explaining and describing at B1.
\end{grammarlens}

\subsection*{Your turn}
\begin{itemize}
\raggedright
  \item \emph{Say it:} \emph{aequitās, aequitātis}
  \item \emph{Say it:} \emph{aequitās, aequitātis}, once more
  \item \emph{Say it:} a sentence of your own with \emph{aequitās, aequitātis}
  \item \emph{Recall:} say the Latin for custody, then the Latin for justice, then say \emph{aequitās, aequitātis} again
\end{itemize}

\begin{tcolorbox}[breakable,colback=teal!4,colframe=teal!35!black,title={Before you move on}]
What does aequitās, aequitātis mean? (\textquotedblleft{}Fairness, equity\textquotedblright{}.) Say it once more.
\end{tcolorbox}
\section[obtemperō, obtemperāre]{obtemperō, obtemperāre — to obey, to comply (with)}

\label{lesson:LA-C416-obtempero}



\begin{quote}
Before the new one: say the Latin for justice, then the Latin for fairness.
\end{quote}

\subsection*{You'll want to know: obtemperō, obtemperāre}
\textbf{obtemperō, obtemperāre} — \textquotedblleft{}to obey, to comply (with)\textquotedblright{}.

\begin{quote}
Lēgibus obtemperāre dēbēmus.
\end{quote}

\begin{quote}
\emph{We must obey the laws.}
\end{quote}

\emph{obtemperō, obtemperāre, obtemperāvī, obtemperātum} + dative, like \emph{pāreō}.

\begin{cousinweb}
\emph{obtemperō} is built from \emph{ob} and \emph{temperō} \textquotedblleft{}to keep within measure, to restrain oneself\textquotedblright{} (EtymDB): to hold oneself in for someone else.
\end{cousinweb}

\begin{grammarlens}[title={What you've built}]
One more word for telling, explaining and describing at B1.
\end{grammarlens}

\subsection*{Your turn}
\begin{itemize}
\raggedright
  \item \emph{Say it:} \emph{obtemperō, obtemperāre}
  \item \emph{Say it:} \emph{obtemperō, obtemperāre}, once more
  \item \emph{Say it:} a sentence of your own with \emph{obtemperō, obtemperāre}
  \item \emph{Recall:} say the Latin for justice, then the Latin for fairness, then say \emph{obtemperō, obtemperāre} again
\end{itemize}

\begin{tcolorbox}[breakable,colback=teal!4,colframe=teal!35!black,title={Before you move on}]
What does obtemperō, obtemperāre mean? (\textquotedblleft{}To obey, to comply (with)\textquotedblright{}.) Say it once more.
\end{tcolorbox}
\section[coerceō, coercēre]{coerceō, coercēre — to restrain, to keep in check}

\label{lesson:LA-C416-coerceo}



\begin{quote}
Before the new one: say the Latin for fairness, then the Latin for to obey.
\end{quote}

\subsection*{You'll want to know: coerceō, coercēre}
\textbf{coerceō, coercēre} — \textquotedblleft{}to restrain, to keep in check\textquotedblright{}.

\begin{quote}
Cōnsul sēditiōsōs coercuit.
\end{quote}

\begin{quote}
\emph{The consul restrained the rebels.}
\end{quote}

\emph{coerceō, coercēre, coercuī, coercitum}, \emph{co-} + \emph{arceō}, to shut in. English \emph{coerce}.

\begin{cousinweb}
\emph{coerceō} is built from \emph{co-} \textquotedblleft{}together\textquotedblright{} and \emph{arceō} \textquotedblleft{}to shut in, to keep off\textquotedblright{} (Lewis \& Short): to hem in on all sides. English \emph{coerce} comes from it.
\end{cousinweb}

\begin{grammarlens}[title={What you've built}]
One more word for telling, explaining and describing at B1.
\end{grammarlens}

\subsection*{Your turn}
\begin{itemize}
\raggedright
  \item \emph{Say it:} \emph{coerceō, coercēre}
  \item \emph{Say it:} \emph{coerceō, coercēre}, once more
  \item \emph{Say it:} a sentence of your own with \emph{coerceō, coercēre}
  \item \emph{Recall:} say the Latin for fairness, then the Latin for to obey, then say \emph{coerceō, coercēre} again
\end{itemize}

\begin{tcolorbox}[breakable,colback=teal!4,colframe=teal!35!black,title={Before you move on}]
What does coerceō, coercēre mean? (\textquotedblleft{}To restrain, to keep in check\textquotedblright{}.) Say it once more.
\end{tcolorbox}
\section[vindicō, vindicāre]{vindicō, vindicāre — to claim; to avenge, to punish}

\label{lesson:LA-C416-vindico}



\begin{quote}
Before the new one: say the Latin for to obey, then the Latin for to restrain.
\end{quote}

\subsection*{You'll want to know: vindicō, vindicāre}
\textbf{vindicō, vindicāre} — \textquotedblleft{}to claim; to avenge, to punish\textquotedblright{}.

\begin{quote}
Iniūriam vindicāvit.
\end{quote}

\begin{quote}
\emph{He avenged the wrong.}
\end{quote}

\emph{vindicō, vindicāre, vindicāvī, vindicātum}. In law, to claim a thing as one's own; \emph{in lībertātem vindicāre}, to set free. English \emph{vindicate}.

\begin{cousinweb}
\emph{vindicō} is built on \emph{vindex, vindicis} \textquotedblleft{}a claimant, a protector, an avenger\textquotedblright{} (EtymDB); Lewis \& Short explain it as \emph{vim dīcō}, to assert force. English \emph{vindicate} comes from it.
\end{cousinweb}

\begin{grammarlens}[title={What you've built}]
One more word for telling, explaining and describing at B1.
\end{grammarlens}

\subsection*{Your turn}
\begin{itemize}
\raggedright
  \item \emph{Say it:} \emph{vindicō, vindicāre}
  \item \emph{Say it:} \emph{vindicō, vindicāre}, once more
  \item \emph{Say it:} a sentence of your own with \emph{vindicō, vindicāre}
  \item \emph{Recall:} say the Latin for to obey, then the Latin for to restrain, then say \emph{vindicō, vindicāre} again
\end{itemize}

\begin{tcolorbox}[breakable,colback=teal!4,colframe=teal!35!black,title={Before you move on}]
What does vindicō, vindicāre mean? (\textquotedblleft{}To claim; to avenge, to punish\textquotedblright{}.) Say it once more.
\end{tcolorbox}
